\chapter{Séries Numéricas}\label{ch:b1:series}

Somar infinitos números significa tomar o limite das somas parciais ---
nada mais, nada menos. Este capítulo estabelece as definições e os
critérios de convergência utilizáveis no primeiro ano: comparação e
equivalentes para termos positivos, o \omterm{thm:b1:series:ratio}{critério da razão}, a comparação
com \omterm{thm:b1:integration:def}{integral} que dá as \omterm{thm:b1:series:riemann}{séries de Riemann}, a \omterm{thm:b1:series:absolute}{convergência absoluta} e o
teorema das \omterm{thm:b1:series:alternating}{séries alternadas}. A teoria mais fina (produtos de \omterm{def:b1:series:def}{séries},
somação por pacotes, \omterm{def:b1:series:def}{séries} de funções) pertence ao segundo ano.

\section{Generalidades}

\begin{definition}\label{def:b1:series:def}
Dada uma sequência $(u_n)$, a \emph{série}\index{série} $\sum u_n$
é a sequência das \emph{somas parciais} $S_N = \sum_{n=0}^{N} u_n$.
A série \emph{converge} quando $(S_N)$ converge; o limite é a
\emph{soma} $\sum_{n=0}^{\infty} u_n$, e $R_N = \sum_{n > N} u_n =
S - S_N$ é o \emph{resto}, que tende a $0$.
\end{definition}

\begin{example}[Série geométrica]\label{ex:b1:series:geometric}
Para $q \in \C$: $\;S_N = \sum_{n=0}^{N} q^n = \frac{1 -
q^{N+1}}{1-q}$ ($q \neq 1$). A \omterm{def:b1:series:def}{série} converge se e somente se $\abs q < 1$
(\cref{exo:b1:seq:3}), com\index{série geométrica}
\[
\sum_{n=0}^{\infty} q^n = \frac{1}{1 - q} .
\]
\end{example}

\begin{example}[Decimais periódicas são séries geométricas]\label{ex:b1:series:periodicdecimal}
Que número é $0.363636\dots$? A sua própria escrita é uma \omterm{def:b1:series:def}{série}:
\[
0.\overline{36} = \sum_{k=1}^{\infty} \frac{36}{100^k}
= 36\cdot\frac{1/100}{1 - 1/100} = \frac{36}{99} =
\frac{4}{11} ,
\]
pela soma geométrica com $q = \frac{1}{100}$. Em geral um
bloco $B$ de $p$ dígitos que se repete indefinidamente vale
$\frac{B}{10^p - 1}$ --- o mecanismo por trás do critério de periodicidade
do \cref{pb:b1:reals:1}, que a linguagem deste capítulo
finalmente enuncia em uma linha: uma \omterm{pb:b1:reals:1}{expansão decimal} é uma \omterm{def:b1:series:def}{série}
convergente, eventualmente periódica exatamente quando a sua soma é
racional. A maquinaria dos dígitos do Capítulo 10, construída lá com
meros \omterm{def:b1:reals:bounds}{supremos}, era a teoria das \omterm{def:b1:series:def}{séries} viajando incógnita.
\end{example}

\begin{proposition}[Primeiros fatos]\label{prop:b1:series:first}
\begin{enumerate}
  \item Se $\sum u_n$ converge, então $u_n \to 0$. (A recíproca é
        \emph{falsa}: a \omterm{def:b1:series:def}{série} harmônica.)
  \item Linearidade: \omterm{def:b1:series:def}{séries} convergentes somam-se e multiplicam-se por
        escalares, com as somas esperadas.
  \item (Telescópica\index{série telescópica}) $\sum (v_{n+1} -
        v_n)$ converge se e
        somente se $(v_n)$ converge, com soma $\lim v_n -
        v_0$.
  \item Alterar um número finito de termos não afeta a convergência
        (apenas a soma).
\end{enumerate}
\end{proposition}

\begin{proof}
(1) $u_N = S_N - S_{N-1} \to S - S = 0$. A \omterm{def:b1:series:def}{série} harmônica tem
$u_n = \frac1n \to 0$ e no entanto diverge (\cref{exo:b1:seq:5}). (2)
Operações sobre limites. (3) $S_N = v_{N+1} - v_0$. (4) As somas
parciais mudam por uma quantidade eventualmente constante.
\end{proof}

\begin{example}[Planejar dígitos com o resto geométrico]\label{ex:b1:series:georemainder}
Para $\abs q < 1$ o resto da \omterm{ex:b1:series:geometric}{série geométrica} é
explícito:
\[
R_N = \sum_{n = N+1}^{\infty} q^n = \frac{q^{N+1}}{1 - q} .
\]
Isso converte metas de precisão em contagens de termos antes de
qualquer cálculo. Para avaliar $\sum_{n\geq0} \bigl(\frac13\bigr)^n =
\frac32$ com erro inferior a $10^{-10}$: é
preciso $\frac{(1/3)^{N+1}}{2/3} \leq
10^{-10}$, isto é, $3^{N} \geq \frac{3}{2}\cdot 10^{10}$, isto é, $N
\geq 22$ (pois $3^{22} \approx 3.1\cdot10^{10}$): vinte e três
termos, conhecidos de antemão. Toda estimativa de taxa geométrica
dos problemas de fim de semana (a \omterm{def:b1:series:def}{série} em $\frac13$ para $\ln 2$, os
\omterm{def:b1:functions:arc}{arco-tangentes} de Machin do \cref{pb:b1:taylor:1}) é este orçamento de duas linhas
em traje profissional.
\end{example}

\begin{example}[Um telescópio mais longo]\label{ex:b1:series:telescope3}
Calcule $\sum_{n\geq1} \frac{1}{n(n+1)(n+2)}$. Frações
parciais (\cref{ch:b1:fractions}):
\[
\frac{1}{n(n+1)(n+2)}
= \frac{1/2}{n} - \frac{1}{n+1} + \frac{1/2}{n+2}
= \frac12\Bigl(\frac{1}{n(n+1)} - \frac{1}{(n+1)(n+2)}\Bigr),
\]
onde a segunda forma --- uma diferença de valores consecutivos de
$w_n = \frac{1}{n(n+1)}$ --- é a telescópica. Logo
\[
\sum_{n=1}^{N} \frac{1}{n(n+1)(n+2)}
= \frac12\Bigl(w_1 - w_{N+1}\Bigr)
= \frac12\Bigl(\frac12 - \frac{1}{(N+1)(N+2)}\Bigr)
\longrightarrow \frac14 .
\]
A ideia de fechamento: frações parciais com três termos raramente
telescopam tal como estão escritas; reagrupe-as primeiro numa
diferença $w_n -
w_{n+1}$ --- a recompensa não é apenas a convergência, mas a
soma exata, que nenhum critério de comparação jamais fornece.
\end{example}

\section{Séries de termos não negativos}

\begin{theorem}[Somas parciais limitadas]\label{thm:b1:series:positive}
Se $u_n \geq 0$ para todo $n$, as somas parciais são crescentes, logo: $\sum
u_n$
converge $\iff$ as suas somas parciais são limitadas superiormente.
Daí o \emph{critério de comparação}: se $0 \leq u_n \leq v_n$ para todo $n$ (grande),
\[
\sum v_n \text{ converge} \implies \sum u_n \text{ converge},
\qquad
\sum u_n \text{ diverge} \implies \sum v_n \text{ diverge}.
\]
E o \emph{critério dos equivalentes}: se $u_n \sim v_n$ com $v_n \geq
0$, as duas
\omterm{def:b1:series:def}{séries} têm a mesma natureza.
\end{theorem}

\begin{proof}
Teorema do limite monótono (\cref{thm:b1:seq:monotone}) para o primeiro
ponto; comparação das somas parciais para o segundo. Equivalentes:
para $n$ grande, $\frac12 v_n \leq u_n \leq 2 v_n$ (definição de $\sim$
com $\varepsilon = \frac12$), e a comparação aplica-se nos dois sentidos.
\end{proof}

\begin{example}[Um equivalente que prova a divergência]\label{ex:b1:series:equivdiverge}
Natureza de $\sum_{n\geq1} n\sin\dfrac{1}{n^2}$? Como
$\frac{1}{n^2} \to 0$ e $\sin h \sim h$ em $0$:
\[
n\sin\frac{1}{n^2} \;\sim\; n\cdot\frac{1}{n^2} = \frac1n ,
\]
e o critério dos equivalentes transfere a divergência da \omterm{def:b1:series:def}{série}
harmônica: divergente --- mesmo que os termos tendam a $0$. Uma
expansão, uma escala, um veredicto; o mesmo padrão em duas etapas
(equivalente e depois consulta às escalas de Riemann ou geométrica)
decide as quatro \omterm{def:b1:series:def}{séries} do \cref{exo:b1:series:3}.
\end{example}

\begin{example}[O critério dos equivalentes em uma linha]\label{ex:b1:series:equivtest}
Natureza de $\sum_{n \geq 1} \frac{\sqrt{n+1} - \sqrt n}{n}$?
Conjugue o numerador:
\[
\frac{\sqrt{n+1} - \sqrt n}{n}
= \frac{1}{n\,(\sqrt{n+1} + \sqrt n)}
\sim \frac{1}{2\,n^{3/2}} ,
\]
uma escala de Riemann convergente ($\alpha = \frac32 > 1$): a \omterm{def:b1:series:def}{série}
converge. Toda a decisão levou um equivalente e uma consulta
--- desde que os termos sejam não negativos, o que é o caso. A
ideia de fechamento: para \omterm{def:b1:series:def}{séries} de termos positivos, toda a teoria
da convergência é um \emph{dicionário de escalas} ($n^{-\alpha}$, $q^n$,
$\frac{1}{n(\ln n)^\alpha}$) mais a licença para substituir um termo
por um equivalente; o trabalho analítico está na assintótica
(\cref{ch:b1:taylor}), nunca na somação.
\end{example}

\begin{theorem}[Comparação com integral; séries de Riemann]\label{thm:b1:series:riemann}
Seja $f$ \omterm{def:b1:continuity:continuous}{contínua}, não negativa e \emph{decrescente} em
$\intco{1}{+\infty}$. Então
\[
\int_1^{N+1} f(t)\,\dd t \;\leq\; \sum_{n=1}^{N} f(n) \;\leq\; f(1)
+ \int_1^{N} f(t)\,\dd t ,
\]
logo $\sum f(n)$ converge se e somente se $\bigl(\int_1^x f\bigr)$ é limitada. Em
particular, para $\alpha \in \R$:\index{séries de Riemann}
\[
\sum_{n \geq 1} \frac{1}{n^\alpha} \text{ converge}
\iff \alpha > 1,
\]
e $\sum_{n=1}^{N} \frac1n = \ln N + O(1)$.
\end{theorem}

\begin{proof}
Para $n \leq t \leq n+1$, a monotonia dá $f(n+1) \leq f(t)
\leq f(n)$; integrando em $\intcc{n}{n+1}$ (um segmento de
comprimento $1$):
\[
f(n+1) \;\leq\; \int_n^{n+1} f(t)\,\dd t \;\leq\; f(n) .
\]
Somando as desigualdades da direita para $n = 1, \dots, N-1$ obtém-se
$\int_1^{N} f \leq \sum_{n=1}^{N-1} f(n)$, donde o enquadramento
superior depois de acrescentar $f(N) \leq f(1)$; somando as da esquerda para
$n = 1, \dots, N$ obtém-se $\sum_{n=2}^{N+1} f(n) \leq
\int_1^{N+1} f$, que após reindexação é o enquadramento inferior.
Convergência: as somas parciais e as integrais $\int_1^x f$
majoram-se mutuamente a menos da constante $f(1)$, e ambas são
não decrescentes, logo uma é limitada se e somente se a outra o é
(\cref{thm:b1:series:positive}). Para $f(t)
= t^{-\alpha}$ ($\alpha \neq 1$): $\int_1^x t^{-\alpha}\dd t =
\frac{x^{1-\alpha} - 1}{1 - \alpha}$, limitada se e somente se $\alpha > 1$;
para $\alpha = 1$ a \omterm{thm:b1:integration:def}{integral} é $\ln x \to \infty$, e o enquadramento
dá $\ln(N+1) \leq H_N \leq 1 + \ln N$. Para $\alpha \leq 0$ os
termos não tendem a $0$.
\end{proof}

\begin{example}[A pilha harmônica]\label{ex:b1:series:harmonicstack}
Quantos termos a \omterm{def:b1:series:def}{série} harmônica precisa acumular para passar de $20$?
O enquadramento $\ln(N+1) \leq H_N \leq 1 + \ln N$ responde sem soma alguma: $H_N \geq 20$ exige $1 + \ln N \geq 20$,
isto é, $N \geq \eu^{19} \approx 1.8\cdot10^{8}$, e está garantido
assim que $\ln(N + 1) \geq 20$, isto é, $N \approx \eu^{20} \approx
4.9\cdot10^{8}$. (O problema de fim de semana afina
isto para $N \approx
\eu^{20 - \gamma} \approx 2.7\cdot10^{8}$ via a \omterm{pb:b1:series:1}{constante de Euler}.)
A ideia de fechamento: a comparação com \omterm{thm:b1:integration:def}{integral} não decide apenas
a convergência --- ela \emph{localiza} as somas parciais com
precisão logarítmica, transformando um cálculo desesperado
(centenas de milhões de termos) numa estimativa de duas linhas.
\end{example}

\begin{theorem}[Critério da razão (d'Alembert)]\label{thm:b1:series:ratio}
Seja $u_n > 0$ com $\frac{u_{n+1}}{u_n} \to \ell$.\index{critério da razão}
\begin{itemize}
  \item Se $\ell < 1$: $\sum u_n$ converge;
  \item se $\ell > 1$: $u_n \to +\infty$, divergência;
  \item se $\ell = 1$: nenhuma conclusão ($\sum \frac1n$ diverge, $\sum
        \frac{1}{n^2}$ converge).
\end{itemize}
\end{theorem}

\begin{proof}
Se $\ell < 1$, fixe $q \in \intoo{\ell}{1}$: a partir de certo $N$,
$u_{n+1} \leq q\,u_n$, logo $u_n \leq u_N q^{\,n-N}$ por indução:
comparação com uma \omterm{ex:b1:series:geometric}{série geométrica}. Se $\ell > 1$: a partir de certo $N$
a sequência $(u_n)$ é crescente, logo não pode tender a $0$
(o seu limite, se existir, é $\geq u_N > 0$); pela
\cref{prop:b1:series:first}~(1), divergência --- e de fato $u_n
\geq u_N q^{n-N}$ com $q > 1$ dá $u_n \to \infty$.
\end{proof}

\begin{example}\label{ex:b1:series:ratio}
$\sum \frac{x^n}{n!}$ converge para todo $x > 0$: razão
$\frac{x}{n+1} \to 0$. A sua soma é $\eu^x$: por Taylor--Lagrange
(\cref{thm:b1:taylor:lagrange}) em $\intcc{0}{x}$,
\[
\Bigl| \eu^x - \sum_{k=0}^{n} \frac{x^k}{k!} \Bigr|
\leq \eu^{x}\, \frac{x^{n+1}}{(n+1)!} \xrightarrow[n\to\infty]{} 0 ,
\]
a cota tendendo a $0$ porque o fatorial domina
(\cref{exo:b1:integration:9}~(1) usou o mesmo fato). O mesmo
argumento soma as \omterm{def:b1:series:def}{séries} de $\sin$, $\cos$, $\sinh$, $\cosh$ em todo
o $\R$.
\end{example}

\begin{example}[O critério da razão é suficiente, não necessário]\label{ex:b1:series:rationotnecessary}
Seja $u_n = 2^{-n}$ para $n$ par e $u_n = 2^{-n-2}$ para $n$ ímpar. As razões
consecutivas oscilam entre $\frac{1}{8}$ e $\frac12\cdot4 = 2$, de modo que $\frac{u_{n+1}}{u_n}$ não tem
limite e d'Alembert é mudo --- e no entanto $u_n \leq 2^{-n}$ e o critério de
comparação resolve a convergência instantaneamente. A hipótese do
critério (a razão \emph{converge}) é uma restrição real: ela serve a
termos com uma única estrutura multiplicativa dominante (fatoriais,
potências), e falha em tudo o que respira. Quando as razões se
comportam mal, recue para a comparação com um envelope geométrico ---
que é tudo o que o \omterm{thm:b1:series:ratio}{critério da razão} sempre foi, como mostra a sua
demonstração.
\end{example}

\begin{example}[O critério da razão em batalhas de fatoriais]\label{ex:b1:series:ratiobinom}
Natureza de $\sum_{n\geq0} \dfrac{(n!)^2}{(2n)!}$ (inversos dos \omterm{def:b1:counting:objects}{coeficientes binomiais} centrais, a
menos do fator $n + 1$)? A razão colapsa os fatoriais:
\[
\frac{u_{n+1}}{u_n}
= \frac{((n+1)!)^2}{(n!)^2}\cdot\frac{(2n)!}{(2n+2)!}
= \frac{(n+1)^2}{(2n+1)(2n+2)}
\longrightarrow \frac14 < 1 :
\]
convergente, com folga --- os termos decaem essencialmente
como $4^{-n}$, coerentemente com $\binom{2n}{n} \geq
\frac{4^n}{2n+1}$ do \cref{pb:b1:integration:1}. A ideia de
fechamento: quocientes de fatoriais são exatamente o que o \omterm{thm:b1:series:ratio}{critério
da razão} digere --- todo fatorial se cancela numa função racional de
$n$, cujo limite se lê nos termos dominantes.
\end{example}

\section{Convergência absoluta; séries alternadas}

\begin{theorem}[Convergência absoluta]\label{thm:b1:series:absolute}
Se $\sum \abs{u_n}$ converge (\emph{convergência
absoluta}\index{convergência absoluta}), então $\sum u_n$ converge,
e $\bigl|\sum u_n\bigr| \leq \sum \abs{u_n}$. Isto vale para termos reais
ou complexos.
\end{theorem}

\begin{proof}
As somas parciais satisfazem, para $M > N$ (critério de Cauchy,
\cref{thm:b1:seq:complete}):
\[
\abs{S_M - S_N} = \Bigl| \sum_{n=N+1}^{M} u_n \Bigr|
\leq \sum_{n=N+1}^{M} \abs{u_n},
\]
que é pequeno para $N$ grande, pois as somas parciais de $\sum\abs{u_n}$
formam uma \omterm{def:b1:seq:cauchy}{sequência de Cauchy}. Logo $(S_N)$ é de Cauchy, portanto
convergente. A desigualdade passa ao limite a partir da desigualdade
triangular finita.
\end{proof}

\begin{example}[Convergência absoluta, real e complexa]\label{ex:b1:series:absexamples}
$\sum_{n\geq1} \frac{\sin n}{n^2}$: os termos mudam de sinal
erraticamente (de fato $(\sin n)$ é \omterm{def:b1:topology:dense}{densa} em $\intcc{-1}{1}$,
\cref{exo:b1:seq:12}), e nenhuma estrutura alternada está à
vista. A \omterm{thm:b1:series:absolute}{convergência absoluta} salva tudo de uma vez:
$\bigl|\frac{\sin n}{n^2}\bigr| \leq \frac{1}{n^2}$, uma
escala convergente, logo a \omterm{def:b1:series:def}{série} converge. O mesmo escudo funciona
sobre $\C$: $\sum_{n\geq1}\frac{\eu^{\iu n}}{n^2}$ converge
porque $\bigl|\frac{\eu^{\iu n}}{n^2}\bigr| = \frac{1}{n^2}$
--- padrões de sinal, mesmo bidimensionais, são irrelevantes
uma vez que os \omterm{def:b1:complex:field}{módulos} sejam somáveis. A ideia de fechamento: a
\omterm{thm:b1:series:absolute}{convergência absoluta} é a única ferramenta deste capítulo que nunca
pergunta como os sinais estão organizados; tente-a primeiro
(\cref{met:b1:series:decide}), e reserve os critérios delicados
para as \omterm{def:b1:series:def}{séries} que a ela resistem.
\end{example}

\begin{theorem}[Critério das séries alternadas]\label{thm:b1:series:alternating}
Seja $(a_n)$ decrescente com $a_n \to 0$. Então a \omterm{def:b1:series:def}{série} alternada
$\sum (-1)^n a_n$ converge; a sua soma está entre duas somas parciais
consecutivas quaisquer, e\index{séries alternadas}
\[
\abs{R_N} = \Bigl| \sum_{n > N} (-1)^n a_n \Bigr| \leq a_{N+1} .
\]
\end{theorem}

\begin{proof}
As somas parciais pares e ímpares são adjacentes: $S_{2p+2} - S_{2p} =
a_{2p+2} - a_{2p+1} \leq 0$
(decrescente), $S_{2p+1} - S_{2p-1} =
a_{2p} - a_{2p+1} \geq 0$ (crescente), e $S_{2p} - S_{2p+1} =
a_{2p+1} \to 0$. Pelo \cref{thm:b1:seq:adjacent} elas partilham
um limite $S$, que o critério das duas \omterm{def:b1:seq:subsequence}{subsequências}
(\cref{prop:b1:seq:subsequences}) torna o limite de $(S_N)$;
além disso $S$ fica preso entre somas parciais consecutivas, e
$\abs{S - S_N}$ é no máximo a distância até a seguinte, $a_{N+1}$.
\end{proof}

\begin{example}[Série harmônica alternada]\label{ex:b1:series:ln2}
$\sum_{n \geq 1} \frac{(-1)^{n-1}}{n}$ converge (critério das alternadas)
mas não absolutamente (\omterm{def:b1:series:def}{série} harmônica). A sua soma é $\ln 2$: da
identidade geométrica finita $\frac{1}{1+t} = \sum_{k=0}^{n-1} (-t)^k +
\frac{(-t)^n}{1+t}$, integre em $\intcc{0}{1}$:
\[
\ln 2 = \sum_{k=1}^{n} \frac{(-1)^{k-1}}{k}
+ (-1)^n \int_0^1 \frac{t^n}{1+t}\,\dd t ,
\qquad
0 \leq \int_0^1 \frac{t^n}{1+t}\,\dd t \leq \frac{1}{n+1} \to 0 .
\]
A convergência é dolorosamente lenta ($R_N \approx \frac{1}{N}$) ---
as \omterm{thm:b1:series:alternating}{séries alternadas} convergem por cancelamento, não por pequenez.
\end{example}

\begin{omfigure}
\begin{tikzpicture}[x=0.78cm, y=7.2cm]
  \draw[-Stealth] (0.3,0.42) -- (11,0.42)
    node[below, font=\small] {$N$};
  \draw (1,0.415) -- (1,0.425) node[below=1pt, font=\small]
    {$1$};
  \draw (5,0.415) -- (5,0.425) node[below=1pt, font=\small]
    {$5$};
  \draw (10,0.415) -- (10,0.425) node[below=1pt, font=\small]
    {$10$};
  \draw[gray, dashed] (0.3,0.6931) -- (10.7,0.6931)
    node[right, font=\small, text=black] {$\ln 2$};
  \node[left, font=\small] at (0.3,1) {$1$};
  \node[left, font=\small] at (0.3,0.5) {$0.5$};
  \draw[omDef, thick]
    plot coordinates {(1,1) (2,0.5) (3,0.8333) (4,0.5833)
    (5,0.7833) (6,0.6167) (7,0.7595) (8,0.6345)
    (9,0.7456) (10,0.6456)};
  \fill[omDef] (1,1) circle (1.5pt);
  \fill[omDef] (2,0.5) circle (1.5pt);
  \fill[omDef] (3,0.8333) circle (1.5pt);
  \fill[omDef] (4,0.5833) circle (1.5pt);
  \fill[omDef] (5,0.7833) circle (1.5pt);
  \fill[omDef] (6,0.6167) circle (1.5pt);
  \fill[omDef] (7,0.7595) circle (1.5pt);
  \fill[omDef] (8,0.6345) circle (1.5pt);
  \fill[omDef] (9,0.7456) circle (1.5pt);
  \fill[omDef] (10,0.6456) circle (1.5pt);
\end{tikzpicture}

{\small As somas parciais $S_N$ da \omterm{def:b1:series:def}{série} harmônica alternada
$1 - \frac12 + \frac13 - \cdots$ saltam por cima do seu limite
$\ln 2$ a cada passo: as somas ímpares por cima, as pares por
baixo, cada salto de tamanho $\frac{1}{N+1}$. O enquadramento é a
demonstração do \cref{thm:b1:series:alternating} tornada visível --- e
o fechamento lento da pinça ($\abs{S_N - \ln 2} \approx
\frac{1}{2N}$, problema de fim de semana
\cref{pb:b1:series:1}) é a razão pela qual ninguém calcula $\ln 2$ assim.}
\end{omfigure}

\begin{remark}[Armadilhas comuns com séries]
(i) \emph{O critério dos equivalentes precisa de um sinal}: sejam
$v_n =
\frac{(-1)^n}{\sqrt n}$ e $u_n = v_n + \frac1n$. Então $\frac{u_n}{v_n} = 1 + \frac{(-1)^n}{\sqrt n} \to 1$, logo $u_n
\sim v_n$; e no entanto $\sum v_n$ converge
(critério das alternadas) enquanto $\sum u_n = \sum v_n + \sum \frac1n$ diverge. A equivalência
controla o \emph{tamanho} dos termos, e para \omterm{def:b1:series:def}{séries} com sinais o
tamanho não é destino --- o critério está enunciado, e é verdadeiro,
apenas para termos (eventualmente) não negativos. (ii) \emph{$u_n \to 0$
não prova nada}: a \omterm{def:b1:series:def}{série} harmônica é o eterno contraexemplo; o
sentido recíproco (\cref{prop:b1:series:first}~(1)) é apenas um teste rápido de
divergência. (iii) \emph{Razão com limite $1$ é silêncio, não
convergência}: tanto $\sum\frac1n$ quanto $\sum\frac{1}{n^2}$ têm razão $\to 1$; passe às
escalas de Riemann ou à comparação com \omterm{thm:b1:integration:def}{integral}. (iv)
\emph{Alternada precisa de decrescente}: $\sum
\frac{(-1)^n}{n + (-1)^n}$ parece alternada e só
é tratada por expansão (\cref{exo:b1:series:5}); o problema de fim de semana do
\cref{ch:b1:taylor} (a questão 23 de lá) mostra que o critério pode falhar por
completo sem monotonia. (v) \emph{Agrupar e reordenar não são de
graça}: inserir parênteses é inofensivo para \omterm{def:b1:series:def}{séries} convergentes mas
pode criar convergência a partir da divergência ($1 - 1 + 1 - \cdots$ agrupado aos
pares), e reordenar pode mudar a própria soma --- o drama encenado
no problema de fim de semana deste capítulo (\cref{pb:b1:series:1}).
\end{remark}

\begin{method}[Decidir a natureza de uma série]\label{met:b1:series:decide}
\begin{enumerate}
  \item Será que $u_n \to 0$? Se não, divergência, pare.
  \item Termos não negativos: procure um equivalente de $u_n$
        (expansões, \cref{ch:b1:taylor}!), compare com escalas de Riemann ou
        geométricas; fatoriais e potências pedem o \omterm{thm:b1:series:ratio}{critério da razão};
        um $f(n)$ decrescente pede a comparação com \omterm{thm:b1:integration:def}{integral}.
  \item Os sinais variam: tente primeiro a \omterm{thm:b1:series:absolute}{convergência absoluta}; se
        falhar, o critério das alternadas (verifique
        \emph{decrescente} com cuidado); além disso, ferramentas do
        segundo ano.
\end{enumerate}
\end{method}

\begin{example}[Denominadores ímpares, meio telescópio]\label{ex:b1:series:oddtelescope}
Calcule $\sum_{n \geq 1} \dfrac{1}{4n^2 - 1}$. Frações
parciais: $\frac{1}{(2n-1)(2n+1)} =
\frac12\bigl(\frac{1}{2n-1} - \frac{1}{2n+1}\bigr)$, logo
\[
\sum_{n=1}^{N} \frac{1}{4n^2 - 1}
= \frac12\Bigl(1 - \frac{1}{2N+1}\Bigr)
\longrightarrow \frac12 .
\]
Compare com $\sum \frac{1}{n(n+1)} = 1$
(\cref{exo:b1:series:1}): mesmo esqueleto telescópico, mas aqui os
termos consecutivos distam dois nos números ímpares, e
o fator $\frac12$ registra o passo. A ideia de fechamento:
telescopar é uma mudança de ponto de vista, não um truque --- sempre
que o termo geral é uma diferença $w_n - w_{n+1}$ de uma sequência
com limite, a soma é $w_1 - \lim w$, exatamente a
\cref{prop:b1:series:first}~(3).
\end{example}

\begin{remark}[O encadeamento da análise, em retrospecto]
Este capítulo é o ponto em que a análise do volume converge, e cada
critério nomeia o seu ancestral. As somas parciais limitadas são o
teorema do limite monótono (\cref{ch:b1:seq}), ele próprio o axioma da
completude do \cref{ch:b1:reals}; a \omterm{thm:b1:series:absolute}{convergência absoluta} é o critério de
Cauchy; o critério \omterm{thm:b1:integration:def}{integral} é o enquadramento de áreas do \cref{ch:b1:integration};
os equivalentes dos termos gerais são as expansões do
\cref{ch:b1:taylor}; e o teorema das alternadas é o lema das \omterm{thm:b1:seq:adjacent}{sequências
adjacentes} em traje domingueiro. Lido de trás para a frente, o
encadeamento explica para \emph{que} servia cada capítulo --- e os
problemas de fim de semana que o atravessam (dígitos $b$-ádicos,
Cesàro--Stolz, as máquinas de irracionalidade, a \omterm{pb:b1:series:1}{constante de Euler})
são as mesmas poucas ideias reencontrando-se em altitude cada vez
maior. A álgebra linear que vem a seguir muda de assunto, não de
padrões: o hábito de enunciados exatos com erro certificado
sobrevive à passagem dos limites às dimensões.
\end{remark}

\begin{remark}[Para onde vão as séries a seguir]
Este capítulo fecha a análise do volume e abre três portas. No
volume do segundo ano de graduação, as \omterm{def:b1:series:def}{séries} adquirem uma variável
($\sum
a_n x^n$: \omterm{def:b1:series:def}{séries} de potências, com o seu raio de convergência) e
depois uma teoria de valores funcionais (\omterm{def:b1:series:def}{séries} de Fourier); a
dicotomia \omterm{thm:b1:series:absolute}{convergência absoluta} versus condicional, dramatizada no
problema de fim de semana abaixo, torna-se a pedra angular de
ambas. Em probabilidade (volume do terceiro ano de graduação), as
esperanças de variáveis aleatórias discretas \emph{são} \omterm{def:b1:series:def}{séries}, e a
\omterm{thm:b1:series:absolute}{convergência absoluta} é o que as torna bem definidas. E a \omterm{def:b1:series:def}{série} de
Riemann $\sum n^{-s}$, levada a $s$ complexo, torna-se a função zeta --- a
\omterm{def:b1:series:def}{série} mais estudada de toda a matemática.
\end{remark}

\section{Exercícios}

\begin{exercise}[$\star$]\label{exo:b1:series:1}
Natureza (e soma, quando telescópica) de:
\[
\sum_{n\geq1} \frac{1}{n(n+1)},
\qquad
\sum_{n\geq2} \ln\Bigl(1 - \frac{1}{n^2}\Bigr),
\qquad
\sum_{n\geq0} \frac{3^n + 4^n}{5^n} .
\]
\end{exercise}

\begin{exercise}[$\star$]\label{exo:b1:series:2}
Natureza de: $\;\sum \dfrac{n^2}{2^n}$; $\;\sum \dfrac{n!}{n^n}$;
$\;\sum \dfrac{2^n\,n!}{n^n}$; $\;\sum \dfrac{3^n\,n!}{n^n}$.
\emph{(\omterm{thm:b1:series:ratio}{Critério da razão}; lembre que $\bigl(1 + \frac1n\bigr)^n \to \eu$.)}
\end{exercise}

\begin{exercise}[$\star$]\label{exo:b1:series:3}
Natureza de: $\;\sum \sin\dfrac{1}{n^2}$; $\;\sum
\Bigl(1 - \cos\dfrac1n\Bigr)$; $\;\sum \dfrac{1}{\sqrt{n(n+1)}}$;
$\;\sum \dfrac{\ln n}{n^2}$ \emph{(compare com $n^{-3/2}$)}.
\end{exercise}

\begin{exercise}[$\star$]\label{exo:b1:series:4}
Demonstre que $\sum_{n\geq1} \frac{1}{n^2}$ converge com soma $\leq 2$,
usando $\frac{1}{n^2} \leq \frac{1}{n(n-1)}$ para $n \geq 2$ e uma majoração
telescópica.
\end{exercise}

\begin{exercise}[$\star\star$]\label{exo:b1:series:5}
Natureza de $\;\sum \dfrac{(-1)^n}{\sqrt n}$, de $\;\sum
\dfrac{(-1)^n}{n + (-1)^n}$ \emph{(expanda: o critério das
alternadas não se aplica diretamente --- por quê?)}, e de $\;\sum
\sin\bigl(\pi\sqrt{n^2+1}\,\bigr)$
\emph{(reduza módulo $\pi$: $\sqrt{n^2+1} = n + \frac{1}{2n} + O(n^{-3})$)}.
\end{exercise}

\begin{exercise}[$\star\star$]\label{exo:b1:series:6}
Para quais $\alpha > 0$ a \omterm{def:b1:series:def}{série} $\sum_{n \geq 2}
\dfrac{1}{n (\ln n)^{\alpha}}$ converge? \emph{(Comparação com
\omterm{thm:b1:integration:def}{integral}; substitua $u = \ln t$.)}
\end{exercise}

\begin{exercise}[$\star\star$]\label{exo:b1:series:7}
Seja $u_n = \dfrac{1}{n} - \ln\Bigl(1 + \dfrac1n\Bigr)$. Demonstre que
$0 \leq u_n \leq \dfrac{1}{2n^2}$, que $\sum u_n$ converge, e
deduza a existência da \omterm{pb:b1:series:1}{constante de Euler}\index{constante de Euler}:
\[
\gamma = \lim_{N \to \infty} \Bigl( \sum_{n=1}^{N} \frac 1n - \ln N
\Bigr) .
\]
\end{exercise}

\begin{exercise}[$\star\star$]\label{exo:b1:series:8}
Calcule as somas
\[
\sum_{n=1}^{\infty} \frac{1}{n(n+2)}
\qquad\text{e}\qquad
\sum_{n=0}^{\infty} \frac{n}{2^n} .
\]
\emph{(Para a primeira: frações parciais. Para a segunda: calcule
$\sum_{n=1}^{N} n x^{n-1}$ em forma \omterm{def:b1:topology:closed}{fechada} e faça $N \to \infty$ em
$x = \frac12$.)}
\end{exercise}

\begin{exercise}[$\star\star\star$]\label{exo:b1:series:9}
(Condensação de Cauchy) Seja $(u_n)$ não negativa e decrescente.
Demonstre que
\[
\sum_{n \geq 1} u_n \text{ converge}
\iff
\sum_{k \geq 0} 2^k\, u_{2^k} \text{ converge},
\]
comparando pacotes de termos entre potências consecutivas de $2$.
Recupere daí o critério de Riemann e o
\cref{exo:b1:series:6}.
\end{exercise}

\begin{exercise}[$\star\star\star$]\label{exo:b1:series:10}
Usando a identidade \omterm{thm:b1:integration:def}{integral} do \cref{ex:b1:series:ln2} adaptada a
$\frac{1}{1+t^2}$, demonstre a fórmula de Leibniz
\[
\frac{\pi}{4} = \sum_{n=0}^{\infty} \frac{(-1)^n}{2n+1}
= 1 - \frac13 + \frac15 - \frac17 + \cdots
\]
com a estimativa de erro $\abs{R_N} \leq \frac{1}{2N+3}$.
\end{exercise}

\begin{exercise}[$\star\star$]\label{exo:b1:series:11}
Natureza de $\displaystyle\sum_{n \geq 1} \frac{1}{n^{1 + 1/n}}$.
\emph{(Calcule o limite de $n^{1/n}$ e encontre um equivalente do
termo geral: o critério de Riemann precisa de um expoente
\emph{fixo}.)}
\end{exercise}

\begin{exercise}[$\star\star\star$]\label{exo:b1:series:12}
Seja $(u_n)$ não negativa e \emph{decrescente} com $\sum u_n$
convergente. Demonstre que $n\,u_n \to 0$ \emph{(majore $n\,u_{2n}$
por uma fatia $\sum_{k=n+1}^{2n} u_k$ e use o critério de Cauchy)}. Mostre que a
recíproca é falsa, e que a hipótese de monotonia não pode ser
removida.
\end{exercise}

\section{Problema: A constante de Euler e a série que muda de
soma}

\begin{problem}[{Problema de fim de semana --- $H_n = \ln n + \gamma +
\frac{1}{2n} + O(n^{-2})$, e rearranjar $1 - \frac12 +
\frac13 - \dots$ para $\frac{\ln 2}{2}$}]
\label{pb:b1:series:1}
Duas histórias partilham a \omterm{def:b1:series:def}{série} harmônica. Primeiro, a
contabilidade exata da sua divergência: $H_n - \ln n$ converge para a
\omterm{pb:b1:series:1}{constante de Euler} $\gamma$ (\cref{exo:b1:series:7}), e este problema
afina o enunciado numa lei bilateral $\frac{1}{2(n+1)}
\leq H_n - \ln n - \gamma \leq \frac{1}{2n}$, certificando $\gamma
= 0.5772\dots$ à mão.
Segundo, o escândalo da convergência condicional: a \omterm{def:b1:series:def}{série} harmônica
alternada soma $\ln 2$ (\cref{ex:b1:series:ln2}), e no entanto \emph{os mesmos termos,
em outra ordem}, somam $\frac{\ln 2}{2}$ --- ou $\ln 2 +
\frac12\ln\frac pq$ para quaisquer $p, q$, ou
qualquer real que se queira (Riemann). As duas histórias são uma
só: as somas rearranjadas são calculadas \emph{com} a lei de
$\gamma$.\index{constante de Euler}\index{rearranjo de séries}

\medskip
\noindent\textbf{Parte I --- $\gamma$, enquadrada.} Ponha $a_n = H_n
- \ln n$ e
$b_n = H_n - \ln(n+1)$.
\begin{enumerate}
  \item Usando $\frac{t}{1+t} \leq \ln(1 + t) \leq t$, mostre que
        $(a_n)$ decresce, $(b_n)$ cresce, e que são adjacentes; o
        seu limite comum é $\gamma$, com $b_n \leq
        \gamma \leq a_n$ para todo $n$.
  \item Primeiro tiro numérico: a partir de $H_{10} = 2.928968\dots$,
        enquadre $\gamma$ entre $b_{10} = 0.5311$ e $a_{10} =
        0.6264$. Que tamanho de $n$
        este enquadramento grosseiro exigiria para quatro decimais?
  \item Mostre a representação exata da cauda $a_n - \gamma =
        \sum_{k \geq n} w_k$ (limite de somas
        parciais), em que
        \[
        w_k = a_k - a_{k+1}
        = \ln\Bigl(1 + \frac1k\Bigr) - \frac{1}{k+1}
        = \int_k^{k+1} \frac{(k + 1 - t)}{t\,(k+1)}\,\dd t ,
        \]
        e deduza da forma integral a majoração bilateral
        $\dfrac{1}{2(k+1)^2} \leq w_k \leq \dfrac{1}{2k(k+1)}$.
\end{enumerate}

\medskip
\noindent\textbf{Parte II --- A lei do $\frac{1}{2n}$.}
\begin{enumerate}[resume]
  \item Some as cotas da questão 3 (ambos os lados telescopam ou
        comparam-se a telescópios) e conclua a lei:
        \[
        \frac{1}{2(n+1)} \;\leq\; H_n - \ln n - \gamma \;\leq\;
        \frac{1}{2n} \qquad (n \geq 1).
        \]
  \item Deduza $H_n = \ln n + \gamma + \frac{1}{2n} +
        O\bigl(\frac{1}{n^2}\bigr)$; mais precisamente, mostre que
        $\gamma_n = H_n - \ln n - \frac{1}{2n}$ satisfaz
        $-\frac{1}{2n(n+1)} \leq \gamma_n - \gamma \leq 0$.
  \item Certifique quatro decimais com $n = 100$: dado $H_{100} =
        5.1873775\dots$,
        calcule $\gamma_{100} = 0.577207\dots$
        e conclua $\gamma = 0.5772 \pm 5\cdot10^{-5}$ (valor
        verdadeiro $0.5772156\dots$).
  \item Dois dividendos da lei, ambos necessários adiante: quando
        $m \to
        \infty$,
        \[
        H_{2m} - H_m = \ln 2 - \frac{1}{4m} +
        O\Bigl(\frac{1}{m^2}\Bigr),
        \qquad
        \sum_{j=1}^{m} \frac{1}{2j-1} = \frac{\ln m}{2} + \ln 2
        + \frac\gamma2 + o(1) ,
        \]
        o segundo via $\sum_{j \leq m} \frac{1}{2j-1} = H_{2m}
        - \frac12 H_m$, e do mesmo modo $\sum_{j=1}^{m}
        \frac{1}{2j} = \frac{\ln m}{2} + \frac\gamma2 + o(1)$.
\end{enumerate}

\medskip
\noindent\textbf{Parte III --- A \omterm{def:b1:series:def}{série} harmônica alternada,
até a segunda ordem.}
\begin{enumerate}[resume]
  \item Mostre (por indução, ou agrupando) a identidade
        $\sum_{k=1}^{2m} \frac{(-1)^{k-1}}{k} = H_{2m} - H_m$,
        e deduza tanto a soma $\ln 2$ (de novo) quanto a velocidade
        exata:
        \[
        \sum_{k=1}^{2m} \frac{(-1)^{k-1}}{k}
        = \ln 2 - \frac{1}{4m} + O\Bigl(\frac{1}{m^2}\Bigr) .
        \]
  \item Deduza o erro assintótico da \omterm{def:b1:series:def}{série} harmônica alternada em
        \emph{qualquer} índice: $S - S_N \sim
        \frac{(-1)^N}{2N}$ --- duas vezes menor do que a
        cota de pior caso $a_{N+1} \approx \frac1N$ do \cref{thm:b1:series:alternating}.
  \item (Aceleração de graça) Mostre que as somas médias
        $\tilde S_N = \frac{S_N + S_{N+1}}{2}$ satisfazem $\tilde
        S_N = \ln 2 + O\bigl(\frac{1}{N^2}\bigr)$. Verifique: $S_{10}
        = 0.64563$, $S_{11} = 0.73654$, $\tilde S_{10} =
        0.69109$,
        contra $\ln 2 = 0.69315$: uma única média compra duas casas decimais.
  \item Explique em duas frases por que nenhum truque desse tipo
        pode ajudar um fenômeno de cauda divergente \emph{positiva}
        como o enquadramento da questão 2: o erro alternado oscila
        (sinal $(-1)^N$), de modo que a média cancela o seu termo
        dominante, enquanto o erro do enquadramento de $\gamma$,
        $\frac{1}{2n}$, tem sinal constante. (Fazer a média de $a_n$ e $b_n$
        \emph{ajuda}: relacione $\frac{a_n + b_n}{2}$ com a estimativa do ponto
        médio $H_n - \ln\bigl(n + \frac12\bigr)$ e mostre que o seu erro é $O\bigl(\frac{1}{n^2}\bigr)$.)
\end{enumerate}

\medskip
\noindent\textbf{Parte IV --- Rigidez e o seu fracasso.}
\begin{enumerate}[resume]
  \item Mostre que a parte positiva $\sum \frac{1}{2j-1}$ e a parte negativa $\sum \frac{1}{2j}$ da
        \omterm{def:b1:series:def}{série} harmônica alternada divergem ambas --- a assinatura da
        convergência \emph{condicional}.
  \item Demonstre o enunciado geral por trás da questão 12: se
        $\sum u_n$ converge mas $\sum \abs{u_n}$ diverge, então a \omterm{def:b1:series:def}{série} das partes
        positivas $\sum u_n^+$ e a das partes negativas $\sum u_n^-$ divergem
        ambas \emph{(a partir de $u_n^\pm = \frac{\abs{u_n} \pm u_n}{2}$: se uma convergisse, a outra
        também convergiria, e portanto $\sum\abs{u_n}$)}.
        Esse reservatório inesgotável de massa positiva e negativa
        é o que a receita de Riemann vai gastar.
  \item (Rigidez) Demonstre: se $\sum u_n$ converge
        \emph{absolutamente} e $\sigma \colon \N \to \N$ é uma bijeção, então $\sum u_{\sigma(n)}$
        converge para a mesma soma \emph{(para $N$ grande os
        primeiros $M$ termos rearranjados contêm $u_0, \dots, u_N$; compare as
        somas parciais através da cauda $\sum_{n > N}\abs{u_n}$)}.
  \item (A receita de Riemann) Seja $t \in \R$. Descreva o rearranjo
        guloso da \omterm{def:b1:series:def}{série} harmônica alternada: tome termos positivos
        $1, \frac13, \frac15, \dots$ até que a soma parcial ultrapasse $t$ pela primeira
        vez, depois termos negativos $-\frac12, -\frac14, \dots$ até que ela caia abaixo
        de $t$ pela primeira vez, e repita. Mostre que cada termo é
        usado exatamente uma vez, que após o primeiro cruzamento as
        somas parciais ficam a uma distância de $t$ menor do que o
        último termo usado, e conclua que a \omterm{def:b1:series:def}{série} rearranjada
        converge para $t$: \emph{qualquer} soma prescrita é
        atingível.
\end{enumerate}

\medskip
\noindent\textbf{Parte V --- A fórmula $(p, q)$.} Fixe inteiros
$p, q \geq 1$. Rearranje a \omterm{def:b1:series:def}{série} harmônica alternada em
blocos: $p$ termos positivos (os próximos inversos ímpares), depois
$q$ termos negativos (os próximos inversos pares), e repita.
\begin{enumerate}[resume]
  \item Verifique que isto é um rearranjo genuíno (cada termo
        exatamente uma vez), e que para $(p, q) = (1, 2)$ ele se lê
        \[
        1 - \frac12 - \frac14 + \frac13 - \frac16 - \frac18 +
        \frac15 - \frac1{10} - \frac1{12} + \dots
        \]
  \item (A metade exata) Para $(p, q) = (1, 2)$, demonstre a
        identidade de blocos
        \[
        \frac{1}{2k-1} - \frac{1}{4k-2} - \frac{1}{4k}
        = \frac12\Bigl(\frac{1}{2k-1} - \frac{1}{2k}\Bigr),
        \]
        e deduza a relação \emph{exata} $T_{3K} = \frac12
        S_{2K}$ entre as somas
        parciais rearranjadas e as originais: a divisão da soma
        pela metade é visível em cada estágio finito, e não apenas
        no limite.
  \item Mostre que a soma parcial após $K$ blocos completos
        vale $\sum_{j=1}^{pK} \frac{1}{2j-1} -
        \sum_{j=1}^{qK} \frac{1}{2j}$, e calcule o seu limite
        com a questão 7:
        \[
        \ln 2 + \frac12 \ln\frac pq .
        \]
  \item Controle as somas parciais \emph{dentro} de um bloco (os
        termos tendem a $0$) e conclua que a \omterm{def:b1:series:def}{série} rearranjada por
        $(p,
        q)$ converge para $\ln 2 + \frac12\ln
        \frac pq$. Em particular $(1, 2)$ dá $\frac{\ln
        2}{2}$:
        verifique contra os nove primeiros termos, $T_9 =
        0.3083$,
        aproximando-se lentamente de $0.3466$.
  \item Verificações de coerência e alcance: $(1,1)$ recupera
        $\ln 2$; $(2,1)$ dá $\frac32\ln 2$; que somas são alcançáveis
        por blocos $(p, q)$, e como este cardápio enumerável se
        compara à carta completa de Riemann (questão 14)?
\end{enumerate}

\medskip
\noindent\textbf{Parte VI --- Epílogo: $\gamma$ em ação, e
síntese.}
\begin{enumerate}[resume]
  \item Identifique a soma da \omterm{def:b1:series:def}{série} convergente
        $\sum_{k\geq1} \bigl(\frac1k - \ln\frac{k+1}{k}\bigr)$
        (\cref{exo:b1:series:7}): mostre que ela vale $\gamma$.
  \item Execute a receita de Riemann (questão 14) para o alvo $t =
        1$
        e liste os doze primeiros termos produzidos ($1, \frac13,
        -\frac12, \frac15, -\frac14, \frac17, \frac19, -\frac16,
        \frac1{11}, \frac1{13}, -\frac18, \frac1{15}$),
        calculando a soma parcial ($\approx 0.980$) --- veja
        o algoritmo respirar em torno do seu alvo.
  \item Mostre que algum rearranjo da \omterm{def:b1:series:def}{série} harmônica alternada
        diverge para $+\infty$ \emph{(blocos de termos positivos
        longos o bastante para ganhar $1$ de cada vez, usando a
        questão 12, separados por termos negativos isolados)}.
  \item Afine o \cref{ex:b1:series:harmonicstack} com a lei de
        $\gamma$: mostre que o primeiro índice com $H_N \geq
        20$
        satisfaz $N = \eu^{\,20 - \gamma}\,(1 + o(1))
        \approx 2.7\cdot10^{8}$ --- a \omterm{pb:b1:series:1}{constante de Euler} é exatamente
        a correção que faltava ao enquadramento grosseiro.
  \item Síntese, uma frase para cada: (i) a lei de $\gamma$ e
        o que cada uma das suas três peças ($\ln n$, $\gamma$,
        $\frac{1}{2n}$) contribui; (ii) por que a convergência
        condicional torna a soma dependente da ordem enquanto a
        \omterm{thm:b1:series:absolute}{convergência absoluta} o proíbe; (iii) como a fórmula
        $(p,q)$ foi um \emph{cálculo} com a lei de $\gamma$, e
        não uma afirmação abstrata de existência; (iv) onde estes
        fios continuam --- \omterm{def:b1:series:def}{séries} de potências e produtos de \omterm{def:b1:series:def}{séries}
        no volume do segundo ano de graduação, e o problema de fim
        de semana do volume do terceiro ano de graduação sobre a
        fórmula de Stirling, em que a mesma contabilidade
        soma-versus-integral funciona a plena potência.
\end{enumerate}
\end{problem}
